\documentclass{article}
\usepackage{amsmath, amssymb, amsthm}

\setlength{\parindent}{0pt}

\newtheorem*{theorem}{Theorem}
\newtheorem*{lemma}{Lemma}
\newtheorem*{corollary}{Corollary}

\begin{document}

\begin{theorem}[Characterisation of constructible numbers]
    A real number is constructible by straightedge and compass if and only if it lies in a field obtained from $\mathbb{Q}$ by 
    a finite sequence of quadratic extensions. In particular, every constructible number is algebraic over $\mathbb{Q}$.
\end{theorem}

\begin{theorem}[Lindemann--Weierstrass (consequence)]
    The number $\pi$ is transcendental.
\end{theorem}

\begin{corollary}
    $\sqrt{\pi}$ is not constructible.
\end{corollary}

\begin{proof}
    If $\sqrt{\pi}$ were constructible, then it would be algebraic over $\mathbb{Q}$. 
    Since algebraic numbers are closed under multiplication,
    \[
    \pi = (\sqrt{\pi})^2
    \]
    would also be algebraic. This contradicts the fact that $\pi$ is transcendental. Hence $\sqrt{\pi}$ is not constructible.
    \end{proof}

    \begin{theorem}
    It is impossible to construct, using straightedge and compass, a square with the same area as a given circle.
    \end{theorem}

    \begin{proof}
    Let a circle of radius $r$ be given. A square of equal area must have side length $s$ satisfying
    \[
    s^2 = \pi r^2 \quad \Rightarrow \quad s = r\sqrt{\pi}.
    \]

    Thus, constructing such a square is equivalent to constructing $\sqrt{\pi}$ from a unit length.

    Since $\sqrt{\pi}$ is not constructible, no such construction is possible.
\end{proof}

\end{document}
